
{\em How should I clean my data? What architecture should I use?}
Training large-scale (i.e., deep) machine learning models entails
making many design decisions.
When making such decisions,
typical practice is to
exhaustively search over a small set of standard options.
For example, we might try a few well-known
data cleaning heuristics,
construct a grid over a hyperparameters,
and choose the options that yield the best models.
However, given that this process explores only a small part of the overall
design space (e.g., one can construct $2^n$ possible training datasets from
a pool of $n$ candidate datapoints),
it is unlikely that this approach really yields the
\textit{optimal} training configuration.

How can we find optimal (or at least, better) training configurations?
To do so,
we take the {optimization} perspective on designing model
training.
From this well-studied
perspective,
deciding on a training configuration---or as we will call it,
a set of {\em metaparameters}---is just a high-dimensional optimization
problem.
The input space of this problem comprises all possible metaparameter choices,
including which datapoints to train on, what model architecture to use,
and how to initialize model weights.
The objective function
takes in a set of metaparameters,
trains a machine learning model according to those metaparameters,
and then returns a target metric evaluated on that model
(e.g., test accuracy).
From this perspective, any procedure for selecting
metaparameters---including the typical practice of
grid-searching over standard options---is just an optimization algorithm,
whose goal is to
maximize the objective function with respect to the (high-dimensional) input.

Given that selecting metaparameters is ``just'' a
high-dimensional optimization problem,
a natural tool to consider is the {\em gradient}.
After all, in many contexts,
gradients offer a more effective approach to maximizing
high-dimensional functions than grid search.
Indeed, for a sufficiently
``well-behaved'' function $\smash{f(x)}$ with gradient $\smash{\nabla f(x)}$,
we can optimize $f$ by
iteratively updating $x$ in the direction of $\nabla f(x)$.
This insight suggests a generic recipe for selecting metaparameters:
first, make the objective differentiable with respect to
the metaparameters; second, update via gradient steps.

Now, the idea of using gradients to search for metaparameters is not new.
Indeed, there is a substantial line of work that aims to optimize metaparameters
(e.g., architectures, regularizers, or data augmentation schemes)
with gradient-based methods
\citep{maclaurin2015gradient,liu2018darts,lorraine2020optimizing}.
However, such methods have not managed to scale beyond relatively
small settings. This state of affairs prompts our main question:
\begin{center}
  \vspace*{-.2em}
  {\em Can we \underline{\smash{scalably}} configure model training
  using gradient-based methods?}
  \vspace*{-.2em}
\end{center}



\begin{figure}
  \centering
  \pgfplotsset{
  compat=1.17,
  cycle list/Set2,
}
\begin{tikzpicture}
  \begin{groupplot}[
      group style={
        group size=3 by 1,
        horizontal sep=1.5cm,
      },
      width=5.5cm, height=5cm,
      /pgf/bar width=1cm,
      ybar,
      xtick=data,
      enlarge x limits=0.5, %
      ymin=0,
      xlabel={Method},
      xmajorgrids=false,    %
      ymajorgrids=true,     %
      grid style={dashed},
      /pgf/bar shift=0cm,
    ]
    \nextgroupplot[
      title={CLIP Data Selection}, 
      symbolic x coords={DataComp \#1, MGD},
      ylabel={$\Delta$ over no filtering (\%)},
      xtick={DataComp \#1, MGD},
    ] 
    \addplot[ybar, fill=gray] coordinates {(DataComp \#1,5)};
    \addplot[ybar, fill=Set2-A] coordinates {(MGD,9)};
      
    \nextgroupplot[
      title={IFT Data Selection}, 
      symbolic x coords={LESS, MGD},
      xtick={LESS, MGD},
      ylabel={$\Delta$ over full dataset (\%)},
    ] 
    \addplot[ybar, fill=gray] coordinates {(LESS,0.25)};
    \addplot[ybar, fill=Set2-A] coordinates {(MGD,1.4)};
      
    \nextgroupplot[
      title={Data Poisoning}, 
      symbolic x coords={GC, MGD},
      xtick={GC, MGD},
      ylabel={Test accuracy drop (\%)},
    ] 
    \addplot[ybar, fill=gray] coordinates {(GC,0.8)};
    \addplot[ybar, fill=Set2-A] coordinates {(MGD,15.4)};
      
  \end{groupplot}
\end{tikzpicture}

  \caption{Our proto-algorithm, metagradient descent (MGD),
    uses gradients to achieve state-of-the-art performance across
    a variety of applications, including data selection and
  data poisoning.}
  \label{fig:headline}
\end{figure}

\subsection{Contributions}
In this work,
we answer this question in the affirmative,
adding ``gradient descent on metaparameters''
to the large-scale machine learning toolkit.
Along the way, we will face---and address---two main challenges.

First, existing methods for computing metagradients do not scale.
In response, we devise an algorithm, $\replay{}$,
that can take
metagradients in large-scale settings.
By combining reverse-mode autodifferentiation (AD)
with an efficient data structure,
\replay{} can calculate exact metagradients for models
with billions of parameters and thousands of training steps.

Second,
we find that metagradients of standard training routines
are not necessarily helpful for optimization,
which we connect to {\em non-smoothness} of the metaparameter
optimization landscape.
Borrowing tools from convex optimization,
we devise a framework for designing
``metasmooth'' training routines that {\em do} admit helpful
metagradients.

Addressing the challenges above
unlocks a simple recipe
for solving a broad range of machine learning tasks:
(a) frame the task as a continuous optimization problem over metaparameters;
(b) design a metasmooth training routine;
(c) perform metagradient descent (MGD).
Applying this recipe:
\begin{itemize}
  \item In the DataComp-\texttt{small} competition \citep{gadre2024datacomp},
  we achieve state-of-the-art pre-training data selection for CLIP
    (2x larger performance improvement than the previous
    DataComp-\texttt{small} leader \citep{ecodatum2024ecodatum});
  \item In the context of data selection for instruction tuning (as introduced by \citet{xia2024less}),
    we substantially improve on data selection for Gemma-2B
    (outperforming existing selection methods
    as well as full-data training);
    \item In the {\em accuracy-degrading} data poisoning setting
    (defined by \citet{huber1964robust} and pioneered by
    \citet{lu2022indiscriminate} for deep neural networks),
      we improve attacks on DNNs by an order of magnitude,
      dropping CIFAR-10 accuracy from $92\%\!\to\! 78\%$
      (the best previous attack \citep{lu2023exploring} 
      only reduces accuracy to $91\%$);
  \item For the task of hyperparameter optimization, we 
    efficiently find a competitive CIFAR-10 learning rate schedule 
    (matching the performance of a schedule found by grid search).
\end{itemize}





\begin{figure}[h]
  \centering
  \begin{tikzpicture}[>=stealth,font=\Large]
    \node (z) at (0,0) {$z$};
    \node (theta) at (3,0) {$\theta = \mathcal{A}(z)$};
    \node (phi) at (6,0) {$\phi(\theta)$};

    \draw[->] (z) -- (theta);
    \draw[->] (theta) -- (phi);

    \draw[blue,dashed,->] (phi) to[bend right=40] node[above]
    {$\nabla_z \phi(\mathcal{A}(z))$} (z);

    \node[above] at (0,-0.9) {\normalsize Training setup};
    \node[above] at (3,-0.9) {\normalsize Trained model};
    \node[above] at (6,-0.9) {\normalsize Observed behavior};
  \end{tikzpicture}
  \caption{An illustration of the metagradient.
    We embed a given aspect of the training setup (e.g., the training
    dataset, or optimizer hyperparameters)
    into a continuous {\em metaparameter} vector $z \in \mathbb{R}^d$.
    This metaparameter defines a model $\mathcal{A}(z)$
    by way of the learning algorithm $\mathcal{A}$,
    which in turn defines an output $\phi(z)$.
    The {\em metagradient}
    $\nabla_z \phi(\mathcal{A}(z)) \in \mathbb{R}^d$
    is the gradient of this model output with respect to
  the metaparameter.}
  \label{fig:meta_gradient}
\end{figure}
